\section{Discussion}
\label{sec:discussion}

\paragraph{No single defense dominates.}
Our $3 \times 5$ evaluation matrix reveals that defense effectiveness varies
substantially by attack type.  Watermark detection reliably identifies
unsigned memory writes but provides no protection if the attacker obtains
the watermark key or bypasses the ingestion path.  \sad{} excels against
\minja{} (bridging-step passages carry high query-aligned similarity)
but is less effective against \injecmem{} (broad-anchor passages are
designed to avoid high similarity to any specific query) and against
\agentpoison{} (the centroid passage targets the mean victim query
embedding, placing it near the boundary of the benign similarity
distribution rather than far above it).
Composite defenses consistently outperform individual defenses across
all three attacks, confirming that no single mechanism covers the full
threat surface and motivating a defense-in-depth portfolio.

\paragraph{Why SAD is difficult against centroid-targeting attacks.}
The centroid passage is crafted to maximise mean cosine similarity across
all victim queries.  In a well-populated memory corpus, many benign entries
share vocabulary with common agent tasks (scheduling, preferences, notes),
so the baseline similarity distribution under calibration already captures
moderate query alignment.  A centroid passage that achieves $\asrr=1.00$
need not be a distributional outlier---it must only exceed the top-$k$
benign entries, not exceed the benign similarity mean by $k$ standard deviations.
This gap between the retrieval criterion and the detection criterion is
a fundamental tension, not an implementation artifact.  Future defenses
may close this gap by calibrating on the top-$k$ retrieved entries rather
than the full corpus mean, tightening the detection boundary.

\paragraph{Trigger prepending is a fundamental protocol requirement.}
The evaluation correction identified in Section~\ref{sec:experiments}
reflects a fundamental asymmetry in the \agentpoison{} threat model:
the attacker controls query content at evaluation time (via a shared
trigger in the user's client), but the defender measures retrieval
without this adversarial query modification.  Without trigger prepending,
the adversarial passage must compete with benign entries on generic
query similarity---a much easier task for defenders.  Future simulations
of trigger-based attacks must implement triggered-query evaluation; we
recommend explicitly documenting the evaluation protocol as part of any
attack benchmark.

\paragraph{Centroid passage approximation.}
Our centroid-targeting passage (Eq.~\ref{eq:centroid}) is a practical
approximation of the optimal adversarial passage that requires full
DPR-based gradient optimization~\citep{chen2024agentpoison}.
At $\asrr=1.00$ in the all-MiniLM-L6-v2 embedding space, this suggests
the centroid approximation is sufficient for this encoder's isotropy
structure.  It may be less effective in embedding spaces with lower
isotropy (e.g., encoder-only models trained on narrow domains), where
a single centroid may not span multiple query clusters.
We leave GPU-accelerated DPR-based trigger optimization for camera-ready.

\paragraph{Generalization across embedding models and memory architectures.}
All experiments use all-MiniLM-L6-v2 with FAISS IndexFlatIP.  Attack
effectiveness and \sad{} calibration will differ for other encoders.
Attacks that rely on trigger-triggered centroid retrieval should transfer
broadly (any dense retriever shares the inner-product objective), but
the absolute \asrr{} values will vary with encoder isotropy and vocabulary.
\sad{}'s calibration is encoder-agnostic in design---it calibrates on
whatever similarity distribution the deployed encoder produces---but its
threshold $k$ may require re-tuning for encoders with different
intra-corpus similarity ranges.  Graph-structured memory systems (e.g.,
those using entity--relation stores) present a different attack surface
that our retrieval-based threat model does not cover.

\paragraph{Deployment recommendations.}
Three practical recommendations emerge from our evaluation.
(i)~\emph{Validate memory provenance at write time}: a watermark-based
check on all incoming memory writes costs $O(|entry|)$ compute and
eliminates unsigned injections entirely.
(ii)~\emph{Monitor per-entry anomaly scores continuously}: a rolling
\sad{} calibration window (100 recent queries, $k=2.0$) adds negligible
retrieval latency while maintaining $\tpr > 0.80$ for bridging-step
attacks.
(iii)~\emph{Apply composite portfolios}: our matrix shows no attack
achieves $\asrr > 0.5$ under both watermark and \sad{} simultaneously
(Table~\ref{tab:attack_defense_matrix}), confirming that layered defenses
provide meaningful compounding protection.

\paragraph{Limitations.}
(i) $\asra$ is modelled rather than measured from live agent execution,
which would require running production LLM agents across thousands of
evaluation queries; modelled values are calibrated to paper-reported
means but do not capture variance from LLM stochasticity.
(ii) The character-level watermark is weaker than the token-level scheme
of \citet{zhao2024watermark}; real-world deployment would embed watermarks
at the LLM token generation layer, giving higher $z$-scores and better
evasion resistance.
(iii) The synthetic corpus of 200 entries may not capture the full
diversity of production memory contents, particularly in multi-user or
multi-domain agent deployments.
(iv) \sad{}'s calibration assumes a stationary query distribution; a
sophisticated adversary who gradually shifts the query distribution over
many sessions could inflate the baseline similarity mean and lower the
effective detection threshold.

\paragraph{Future work.}
Key extensions include: (a)~GPU-accelerated DPR-based trigger optimization
matching the original paper's gradient procedure~\citep{chen2024agentpoison};
(b)~live agent integration (LangChain, OpenAI Assistants API) for
direct $\asra$ measurement; (c)~token-level unigram watermark via an LLM
backbone for full paper-fidelity comparison with \citet{zhao2024watermark};
(d)~multi-agent propagation experiments where a single poisoned entry
spreads across a shared knowledge base serving multiple agents;
(e)~formal robustness analysis for composite portfolios under
simultaneously-adaptive adversaries; and (f)~evaluation on
\citet{memorygraft2025} and \citet{amemguard2025} architectures
to assess cross-system generalization.
